\documentclass[noauthor]{ximera}

\title{Week 4}
\author{Math 1193 Design Team}

\begin{document}
\begin{abstract}
    Solutions for Week 4.
\end{abstract}
\maketitle

\begin{problem}
Fill in as much of the following table as you can, then 
compare with your group. Are there any cells you cannot fill in?

\begin{center}
\begin{tabular}{|c|c|c|c|} 
 \hline
 \(x\) & \(f(x)\) & \(3f(x)\) & \(f(3x)\) \\
 \hline\hline
 -3 & 0 & 0 &  \\ 
 \hline
 -1 & -2 & -6 & 0 \\
 \hline
 0 & 7 & 21 & 7 \\
 \hline
 1 & 3 & 9 & -1 \\
 \hline
 3 & -1 & -3 &  \\
 \hline
\end{tabular}
\end{center}

\begin{explanation}
For the \(3f(x)\) column, we're multiplying all \(f(x)\) values by 3.

For the \(f(3x)\) column, we're calculating \(3x\), then seeing which value in the second column corresponds to that input. Notice that a few cells are blank,
since \(f(-9)\) and \(f(9)\) are not defined in the table.
\end{explanation}
\end{problem}

\begin{problem}
Fill in as much of the following table as you can, then 
compare with your group. Are there any cells you cannot fill in?

\begin{center}
\begin{tabular}{|c|c|c|c|} 
 \hline
 \(x\) & \(f(x)\) & \(f(x)-1\) & \(f(x - 1)\) \\
 \hline\hline
 1 & 2 & 1 &  \\ 
 \hline
 2 & 6 & 5 & 2 \\
 \hline
 3 & 4 & 3 & 6 \\
 \hline
 4 & 9 & 8 & 4 \\
 \hline
 5 & -1 & -2 & 9 \\
 \hline
\end{tabular}
\end{center}

\begin{explanation}
    For the \(f(x)-1\) column, we're subtracting 1 from all \(f(x)\) values.

For the \(f(x - 1)\) column, we're calculating \(x - 1\), then seeing which value in the second column corresponds to that input. Notice that one cell
is blank, since \(f(0)\) is not defined in the table.
\end{explanation}
\end{problem}

\begin{problem}
Fill in as much of the following table as you can, then 
compare with your group. Are there any cells you cannot fill in?

\begin{center}
\begin{tabular}{|c|c|c|c|} 
 \hline
 \(x\) & \(f(x)\) & \(-f(x)\) & \(f(-x)\) \\
 \hline\hline
 -2 & -2 & 2 & 6 \\ 
 \hline
 -1 & 0 & 0 & 4 \\
 \hline
 0 & 2 & -2 & 2 \\
 \hline
 1 & 4 & -4 & 0 \\
 \hline
 2 & 6 & -6 & -2 \\
 \hline
\end{tabular}
\end{center}

\begin{explanation}
    For the \(-f(x)\) column, we're multiplying all \(f(x)\) values by \(-1\).

For the \(f(-x)\) column, we're calculating \(-x\), then seeing which value in the second column corresponds to that input.
\end{explanation}

\end{problem}

\begin{problem}
The graph of function \(g\) is a vertical translation of 
the graph of function \(f\).

\begin{center}
    \includegraphics{transformationsVisual.png}
\end{center}

\begin{enumerate}
    \item Complete the \(g(x)\) column of the table.
    \textbf{Do not fill out the last column yet.}
    \begin{center}
        \begin{tabular}{|c|c|c|c|} 
        \hline
        \(x\) & \(f(x)\) & \(g(x)\) & \(h(x) = f(x) - 2.5\) \\
        \hline\hline
        -4 & 0 & 4 & -2.5 \\ 
        \hline
        -3 & -5.8 & -1.8 & -8.3 \\
        \hline
        -0.7 & 0 & 4 & -2.5 \\
        \hline
        1.2 & -3.3 & 0.7 & -5.8 \\
        \hline
        2 & 0 & 4 & -2.5 \\
        \hline
        \end{tabular}
    \end{center}
    \item If \(f(0) = -0.86\), what is \(g(0)\)? Explain to 
    a partner how you know.
    \item Write an equation for \(g(x)\) in terms of \(f(x)\) for any
    input \(x\).
    \item The function \(h\) can be written in terms of \(f\) as 
    \(h(x) = f(x) - 2.5\). Complete the \(h(x)\) column of the 
    table.
\end{enumerate}

\begin{explanation}
    \begin{enumerate}
        \item Since \(f(-4) = 0\) and \(g(-4) = 4\), we know that the graph is shifted up by 4 units. This upward shift corresponds to adding 4 to each \(y\)-coordinate.
        \item Since every \(y\)-coordinate on the graph of \(g\) is 4 more than the \(y\)-coordinate on the graph of \(f\), \(g(0) = f(0) + 4 = -0.86 + 4 = 3.14\).
        \item \(g(x) = f(x) + 4\)
        \item Here, we subtract 2.5 from each \(f(x)\) value.
    \end{enumerate}
\end{explanation}
\end{problem}

\begin{problem}
In this exercise, we'll draw a graph whose formula is 
difficult to write down concisely.

\begin{enumerate}
    \item Draw the line \(y = 2x - 1\), but only for
    \(-5 < x \le -1\).
    \item On the same coordinate plane, draw the line 
    \(y = -x - 1\), but only for \(-1 < x < 2\).
    \item On the same coordinate plane, draw the line \(y = x\),
    but only for \(2 \le x \le 4\).
    \item Is the resulting graph the graph of a function? Why or
    why not? Discuss with a partner.
\end{enumerate}

\begin{explanation}
One way to do this is to graph each line, then erase the points whose \(x\)-coordinates don't satisfy the inequality. Remember to use open and closed dots to represent 
the function's actual value at \(-5\), \(-1\), \(2\), and \(4\).

\begin{center}
    \includegraphics[width=10cm]{piecewise.png}
    \end{center}
Since the graph passes the vertical line test (each input has only one output/each \(x\)-value has only one \(y\)-value), this the graph of a function.
\end{explanation}
\end{problem}


\end{document}
